% a06_BALtorus_d2b: continuation of a06_BALtorus_d2 (input from sec/a06_BALtorus_c.tex): grouped families and the geometry of the regions.

\begin{lemma}[grouped families]\label{lem:a06d:K2}
Let $X$ be any of $\Delta_k[\mathfrak D]$, $\mathcal R_k[\mathfrak D]-\mathcal R_k[\mathfrak D']$, $\mathrm{out}_k[\mathfrak D;C]-\mathrm{out}_k[\mathfrak D;C']$, written with the line telescoping of \cref{lem:a06d:V5} (order induced on $G/S$ from $G$), the symmetric expansion of \cref{lem:a06d:V7}(b) for vertex multipliers and sites, and the multilinear expansion in contraction kernels. Then $X$ is a finite sum, of at most $C(\bar n)$ terms, of renormalised families in the sense of \cite[p.~420]{Ba3} (each divergent two-scalar-leg subgraph $S$ accompanied by $G/S$ with the \cite[(1.7)]{Ba3} vertex and the same changed lines, legs and sites outside $S$), each with exactly the designated object of \cref{lem:a06d:V9}(b)--(d).
\end{lemma}

\begin{proof}
A term with $\delta_k$ on a line of $S$ has $S$ convergent ($\delta_k$ is regular on the diagonal); a term with $\delta_k$ off $S$ pairs with the $G/S$ term with $\delta_k$ on the corresponding line, and $S-c_S=[S-S(T)]+[S(T)-c_S]$, both convergent ($c_S$ uses torus lines, \cref{conv:a06d:LOC}(b), \cite[(1.31)]{Ba3}). The vertex expansion is symmetric, so the subsets of changed legs and sites of $S$ and of its partner coincide. Contraction kernels are lifted index-$k$ kernels (regular) or their deep differences. Polarisation subgraphs keep bubble and seagull with the same multipliers on corresponding legs, and the local cancellation survives $C^2$ multipliers, the subtraction $\sum_{x'}\Pi(x,x')[m(x')-m(x)-\nabla m(x)\cdot(x'-x)]$ being UV-finite in $d=3$. The grouping is exactly \cite[p.~414]{Ba3}'s line list $\{G(\Omega),G(\Omega_2),\delta G\}$.
\end{proof}

\begin{lemma}[geometry]\label{lem:a06d:V8}
\textup{(a)} $\dist(\Lambda_{i+1},\Lambda_i^c)>r$ for every $i$ (the definition \cite[(2.8)]{Ba2}, with $r=r(t_k)$ at scale $k$). \textup{(b)} At scale $k$, with $c_d$ a dimensional constant: $\dist(B^k(\Lambda^{(k)}_7),B^k(\Lambda^{(k)}_5)^c)>2r(t_k)-c_d$, $\dist(B^k(\Lambda^{(k)}_7),B^k(\Lambda^{(k)}_2)^c)>5r(t_k)-c_d$, and $\dist(B^k(\Lambda^{(k)}_7),B^k(\Lambda^{(k-1)'}_7)^c)>7r(t_k)-c_d$, because $\Lambda^{(k)}_7\subset\Lambda^{(k)}_0\subset\Lambda^{(k-1)'}_7$ \cite[p.~570]{Ba2}. \textup{(c)} $\rho(\mathfrak D_k)\ge(5/L)r(t_{k-1})-c_d\ge r(t_k)-c_d$ at $L=5$ (the primed sets are the scale-$(k-1)$ sets read on $T^{(k)}_1$, distances divided by $L$, inclusions kept). \textup{(d)} If $7r(t_k)-c_d\ge M$ then $g_k\equiv1$ on $B^k(\Lambda^{(k)}_7)$, and the coefficient \cite[(1.29)]{Ba3} there differs from the full $\dm$ only by tails $e^{-\delta(7r-M)}$. \textup{(e)} Put $\mathrm{Col}_k\defeq\Lambda^{(k-1)'}_7\cap\Lambda^{(k)c}_7$ and let $D<7r(t_k)-c_d$; every scale-$k$ block $b\subset\Lambda^{(k)}_7$ with $\dist(b,\Lambda^{(k)c}_7)\le D$ lies within $D+1$ of a block of $\mathrm{Col}_k$, hence for $\delta>0$, $\sum_{b\subset\Lambda^{(k)}_7}e^{-\delta\dist(b,\Lambda^{(k)c}_7)}\le C(\delta)|\mathrm{Col}_k|_k+e^{-\delta(7r(t_k)-c_d)}N_k$ with $C(\delta)\defeq\sum_{D\ge0}(2D+3)^3e^{-\delta D}$ and $N_k\defeq|\Lambda^{(k)}_7|_k$. \textup{(f)} Any zone $Z$ of blocks within distance $D_Z<7r(t_k)-c_d$ of $\partial B^k(\Lambda^{(k)}_7)$, on either side, has $|Z|_k\le(1+(2D_Z+3)^3)|\mathrm{Col}_k|_k$.
\end{lemma}

\begin{proof}
(a): by \cite[(2.8)]{Ba2}, $\Lambda^c_{i+1}$ is the union of the large blocks at distance $\le r$ from $\Lambda^c_i$. (b): iterate (a) along a shortest path; each crossing of a shell costs at least $r$ up to the block size. (c), (d): from (b) and the definition of $g_k$. (e): the nearest point of $\Lambda^{(k)c}_7$ to $b$ is within $D$, and since $(\Lambda^{(k-1)'}_7)^c$ is farther than $7r-c_d>D$ that point lies in $\mathrm{Col}_k$; the map $b\mapsto$ (a collar block within $D+1$) has fibres of size $\le(2D+3)^3$; sum over $D$. (f): the blocks of $Z$ outside $\Lambda^{(k)}_7$ lie in $\mathrm{Col}_k$ (at most $|\mathrm{Col}_k|_k$ of them), those inside are covered by (e).
\end{proof}
